\documentclass[10pt]{letter}
\usepackage[margin=0.65in]{geometry}
\usepackage{hyperref}
\usepackage{url}
\setlength{\parskip}{3pt}
\setlength{\parindent}{0pt}
\setlength{\topskip}{4pt}
\addtolength{\topmargin}{-0.3in}
\addtolength{\textheight}{100pt}

\signature{Lizhuo Zhang \\ zhanglizhuo@hunau.edu.cn}
\address{Lizhuo Zhang \\ Hunan Agricultural University \\ Changsha 410128, China}

\begin{document}

\begin{letter}{Editor-in-Chief \\ Neurocomputing}

\opening{Dear Editor-in-Chief,}

We submit the enclosed manuscript, ``Shared Class-Neighborhood Geometry: A Candidate Mechanism for Why CLIP Predicts MLLM Difficulty,'' as a Research Paper for consideration in \emph{Neurocomputing}. Lizhuo Zhang (zhanglizhuo@hunau.edu.cn) is the corresponding author; Yan Ma has approved this submission.

The question we address is fundamental to understanding learning systems: why does the per-class zero-shot accuracy of CLIP, a frozen 2021 model with a simple contrastive objective, predict the per-class accuracy of multimodal large language models (MLLMs) that are one to two orders of magnitude larger and trained with autoregressive language modeling? We test a specific mechanism (shared class-neighborhood geometry) through five lines of evidence: (i) the correlation survives replacing CLIP's encoder with non-CLIP encoders (SigLIP $\rho=0.716$, QwenViT $\rho=0.669$; twelve MLLMs, four encoder families), ruling out the shared-encoder account; (ii) class-centroid geometry agrees strongly across seven independently trained encoders, including a supervised ImageNet-1k-only ResNet-50, replicating from 8 to 40 to 102 classes (Mantel $\rho=0.51$--$0.98$); (iii) the shared structure localizes to class prototypes rather than per-patch features (adapter $R^2\le0.21$ vs.\ Mantel $\rho\ge0.90$); (iv) bootstrap mediation attributes a substantial share of the CLIP--MLLM correlation to shared per-class difficulty on Stanford40 (all six indirect-effect CIs exclude zero), with counterfactual mediator controls attributing nothing to permuted or noise mediators; and (v) CLIP's most dramatic per-class failures are repaired by prompt engineering alone, across all five low-accuracy classes ($0\%\rightarrow100\%$ for ``answer''), locating them in text--image alignment rather than absent visual information.

We are deliberately explicit about what this paper does \emph{not} claim: all analyses are observational, and the geometry$\rightarrow$difficulty link is correlational. Two MLLM perturbation attempts were inconclusive and are reported in full, including their failure modes. A preregistered encoder intervention (frozen language model, swapped encoders, equal-budget projectors) is likewise reported as an honest boundary: every trained arm acquires genuine discriminative power (36.1\%/30.0\%/35.2\% vs.\ 8.4\% random-encoder control), yet cross-arm difficulty profiles do not track geometry agreement above the shared-LM floor. Standard robustness checks accompany every central claim (class-frequency partial control, jackknife leverage, FDR-corrected perturbation grid, complete low-accuracy coverage in the prompt experiment). The contribution is the measured, falsification-tested structural bridge from cross-encoder geometric convergence to a specific cross-model behavioral regularity, not a causal proof.

This manuscript is complementary to, and explicitly differentiated from, our prior papers. Our companion paper in \emph{IEEE Access} (Ma \& Zhang, 2026) established the predictive utility of AnchorScore as a per-class difficulty diagnostic; the present work investigates the representation-level structure that may explain its transferability across heterogeneous vision encoders, and the incremental contribution is the mechanism, not the predictor. A parallel manuscript under consideration elsewhere addresses annotation reliability workflows and shares no tables or figures with this paper.

We believe this work fits \emph{Neurocomputing} on all three of its axes: theory, practice, and applications. Theoretically, it advances the understanding of learning systems by localizing where independently trained large-scale vision encoders converge (class-prototype geometry) and where they do not (per-patch features), supported by exact permutation Mantel tests, bootstrap mediation with counterfactual controls, and explicit power analysis. Practically, it explains why a minutes-computable diagnostic (AnchorScore) transfers across heterogeneous vision encoders, with computational validation across three domains, seven vision encoders, and twelve MLLMs. In application, it supplies a validated lightweight tool for per-class difficulty prediction in MLLM evaluation, together with a boundary result from a preregistered encoder intervention that delimits when the geometry-to-behavior link does and does not hold. All scripts and results are public at \url{https://github.com/zhanglizhuo/SharedVisual}.

This manuscript is original, has not been published elsewhere, and is not under consideration by another journal. All authors have read and approved the final version. The authors declare no competing interests, and this research received no external funding. SCB5 is public on Hugging Face; code, prompts, and result files are open at \url{https://github.com/zhanglizhuo/SharedVisual}.

Thank you for your consideration.

\closing{Sincerely,}

\end{letter}

\end{document}
